\begin{theorem}[Common physical density with geometrically derived exterior walks]%
    \label{thm:peps_torus_simply_connected_physical_density}
    \lean{TNLean.PEPS.IsGIsometric.exists_torusPhysicalCut_common_density_of_isSimplyConnected}
    \leanok
    \uses{def:peps_renyi_entropy}
    Work on a native torus with both dimensions at least three and
    an original regular $G$-isometric site map. Assume the actual
    closed-cell realization $X_R$ is simply connected and the induced
    occupied graph is connected. Let $b$ be the number of crossing
    bonds. There exists one physical reduced density
    $\rho_R$ for every nonzero coherent superposition of commuting
    torus closures, with coefficient matrix $M$. It satisfies
    \begin{align}
        \frac{MM^\dagger}{\operatorname{tr}(MM^\dagger)} & =\rho_R,             &
        \rho_R                                           & \geq0,               & \operatorname{tr}\rho_R & =1,
        \label{eq:peps_simply_connected_density}                                                                \\
        \operatorname{rank}\rho_R                        & =|G|^{b-1},          &
        \rho_R^2                                         & =|G|^{-(b-1)}\rho_R,
        \label{eq:peps_simply_connected_density_flat}                                                           \\
        S(\rho_R)                                        & =(b-1)\log|G|,       &
        S_\alpha(\rho_R)                                 & =(b-1)\log|G|.
        \label{eq:peps_simply_connected_density_entropy}
    \end{align}
    The R\'enyi formula holds for every nonnegative finite order,
    including zero and the von Neumann value at one.
    The exterior walks, spanning trees, roots, and boundary enumeration
    are derived internally. This is a specialization of
    \cite[Theorem~6.9]{Schuch2010PEPS} to regular native tori with
    periods at least three. Occupied graph connectivity expresses
    the connected block in the source, lines 1931--1933 and
    1349--1354; simple connectedness of the closed-cell union alone
    does not force four-neighbor connectivity of the occupied sites.
\end{theorem}

\begin{proof}\leanok
    \uses{thm:peps_torus_physical_cut_density,
        thm:peps_torus_simply_connected_boundary_loops,
        thm:peps_torus_simply_connected_complement}
    The induced complement graph is connected. Choose spanning trees
    and roots in both regions and enumerate the nonempty boundary.
    The simply connected planar lift and its exterior collar provide
    complementary boundary loops with the winding of the two-tree
    boundary routes. Applying
    \ref{thm:peps_torus_physical_cut_density} to these derived loops
    gives (\ref{eq:peps_simply_connected_density}),
    (\ref{eq:peps_simply_connected_density_flat}), and
    (\ref{eq:peps_simply_connected_density_entropy}).
\end{proof}

\begin{theorem}[Physical local equivalence with geometrically derived exterior walks]%
    \label{thm:peps_torus_simply_connected_physical_local_equivalence}
    \lean{TNLean.PEPS.IsGIsometric.torusPhysicalCut_local_equivalence_of_isSimplyConnected}
    \leanok
    \uses{thm:peps_torus_simply_connected_physical_density,
        def:peps_normalized_pure_cut_matrix}
    Under the hypotheses of
    \ref{thm:peps_torus_simply_connected_physical_density}, let $A$ and
    $B$ be the physical coefficient matrices of any two nonzero
    coherent superpositions of commuting torus closures.
    Every operator $O$ on the occupied physical region has equal
    normalized expectation in these states. There exists a unitary
    $U$ on the complementary physical factor satisfying
    \begin{align}
        \widehat B & =\widehat A U^{\mathsf T},                         &
        \operatorname{tr}(O\widehat B\widehat B^\dagger)
                   & =\operatorname{tr}(O\widehat A\widehat A^\dagger).
        \label{eq:peps_simply_connected_local_equivalence}
    \end{align}
    This is a specialization of
    \cite[Theorem~6.7 and Corollary~6.8]{Schuch2010PEPS} to connected
    simply connected regions of regular native tori with periods at
    least three. The width-one stripe construction is separate.
\end{theorem}

\begin{proof}\leanok
    \uses{thm:peps_torus_physical_local_equivalence,
        thm:peps_torus_simply_connected_boundary_loops,
        thm:peps_torus_simply_connected_complement}
    Derive the complementary winding loops from the actual simple
    connectedness and apply
    \ref{thm:peps_torus_physical_local_equivalence}.
\end{proof}
